../cictikz/src/cictikz/data/tex/cictikz_preamble.tex

% LDO with an NMOS pass-fet. The 1.2 V reference is divided down to the
% positive OTA input, the output feeds back to the negative input, and
% the OTA drives the NMOS gate. The source follower does not invert, so
% feedback to the negative input closes a negative loop.

\begin{circuitikz}[american, thick, transform shape, circuit ee IEC]

  ../cictikz/src/cictikz/data/tex/cictikz_lib.tex

  % reference divider
  \draw (0,6.0) to [short,o-] (0,6.0);
  \node[anchor=south] at (0,6.15) {$1.2 V$};
  \draw (0,4.4) \vresistor{};
  \draw (0,2.8) \vresistor{};
  \draw (0,2.8) \vground;
  \fill (0,4.4) circle (0.075);

  % OTA, plus on top from the divider, minus from the feedback
  \draw (0,4.4) -- (1.6,4.4);
  \draw (1.6,4.4) \cicOtaSSWP{$+$}{$-$};

  % NMOS pass-fet, source follower from the 1.5 V input
  \draw (6.4,3.2) \lvnmos{M1}{};
  \draw (6.4,4.8) \vsupply;
  \node[anchor=south] at (6.4,5.25) {$1.5 V$};
  \draw (cicOtaS_out) -- (M1.G);

  % output node and load
  \fill (6.4,3.2) circle (0.075);
  \draw (6.4,3.2) -- (7.8,3.2);
  \node[anchor=south] at (7.8,3.3) {$0.8 V$};
  \draw (7.8,3.2) to [vR, l=$I_{load}$] (7.8,1.6);
  \draw (7.8,1.6) \vground;

  % feedback to the negative input
  \draw (6.4,3.2) -- (6.4,2.4) -- (1.2,2.4) -- (1.2,3.6) -- (1.6,3.6);

\end{circuitikz}
\end{document}
